\subsection*{\texttt{\_\_device\_\_} 与普通 C++ 内联代码}

\co@__device__@ 函数与普通 C++ 函数都会被保留原样，但二者进入不同的编译路径——这个区别很重要。

\co@__device__@ 函数编译成设备代码；\co@__co__@ 函数体用 \co@call@ 关键字调用它，生成的 \co@__global__@ 核函数会在设备上直接调用。

普通 C++ 函数（没有 \co@__device__@ 或 \co@__co__@ 标注）编译成 \textbf{CPU 代码}，运行在主机上，负责建立缓冲区、启动核函数、打印结果。不能在 \co@__co__@ 函数体里调用主机函数。

边界很严格：\co@__device__@ 在 \co@c++@ 划分的设备侧，主机 C++ 在主机侧。\co@__co__@ 函数是桥——编译器把它变换成 \co@__global__@ 核函数（设备侧），同时生成启动它的主机侧包装。
